% Part: counterfactuals
% Chapter: minimal-change-semantics
% Section: antecedent-strengthening

\documentclass[../../../include/open-logic-section]{subfiles}

\begin{document}

\olfileid[tr]{cnt}{min}{agg}

\olsection{Önbileşeni Güçlendirme}

``Önbileşeni güçlendirme'', $!A\lif !C\Entails(!A\land !B)\lif !C$
çıkarımını ifade eder. Maddi koşul için geçerli, karşıolgusallar için
geçersizdir. Bu kibriti çaksaydım yanacağının doğru olduğunu varsayalım.
(Yani kibritte ya da kibrit kutusunun yüzeyinde bir sorun yoktur, kibriti
kırmayacağımdır vb.) Fakat bu kibriti uzay boşluğunda çaksaydım yanacağı
doğru değildir. Dolayısıyla şu çıkarım geçersizdir:
\begin{quote}
  Kibrit çakılsaydı yanardı.

  Öyleyse kibrit uzay boşluğunda çakılsaydı yanardı.
\end{quote}

Lewis--Stalnaker koşul açıklaması bunu açıklar: Kibriti uzay boşluğunda
çaktığım en yakın dünya, yalnızca kibriti çaktığım en yakın dünyaya göre
gerçek dünyadan çok daha uzaktır. Bu nedenle kibrit ikinci dünyada yansa da
birincisinde yanmaz. Olması gereken de budur.

\begin{ex}
  Küre semantiği bu çıkarımı geçersiz kılar; yani
  $p\cif r\Entails/(p\land q)\cif r$. $W=\{w,w_1,w_2\}$,
  $O_w=\{\{w\},\{w,w_1\},\{w,w_1,w_2\}\}$,
  $V(p)=\{w_1,w_2\}$, $V(q)=\{w_2\}$ ve $V(r)=\{w_1\}$ olan
  $\mModel{M}=\tuple{W,O,V}$ modelini ele alalım. $p$'yi kabul eden
  $S=\{w,w_1\}$ küresi vardır ve $p\lif r$ bu kürenin bütün dünyalarında
  doğrudur; dolayısıyla $\mSat{M}{p\cif r}[w]$. Ayrıca $p\land q$'yu kabul
  eden $S'=\{w,w_1,w_2\}$ küresi vardır, fakat
  $\mSat/{M}{(p\land q)\lif r}[w_2]$; dolayısıyla
  $\mSat/{M}{(p\land q)\cif r}[w]$
  (bkz. \olref{fig:antecedent-strengthening}).

\begin{figure}
\begin{center}
\begin{tikzpicture}[layerwidth=1.5,scale=.6]\tiny
  \spheresystem{3}
  \propositionintersect{-10}{3}{30}{5}
  \begin{scope}
    \clip plot[smooth,tension=1] coordinates
          {(0,4.5) (30:.95) (4.5,-3.5)} -- (4.5,4.5) ;
    \spherefill{2}
  \end{scope}
  \draw[proposition] plot[dashed,smooth,tension=1] coordinates
          {(0,4.5) (30:.95) (4.5,-3.5)} ;
  \proposition{30}{2}{30}{5}
  \path (0,0) node[world] {$w$};
  \spherepos{30}{2}{node[world] {$w_1$}}
  \spherepos{-10}{3}{node[world] {$w_2$}}
  \spherepos{-10}{4}{node {$q$}}
  \spherepos{30}{4}{node {$r$}}
  \draw (1,4.5) node[below] {$p$};
\end{tikzpicture}
\caption{Önbileşeni güçlendirmeye karşı örnek}
\ollabel{fig:antecedent-strengthening}
\end{center}
\end{figure}
\end{ex}

\end{document}
